% =========================================================
% Notes: Number Systems --- Axioms, Properties, and Comparison
% File: volume-ii/numbers-summary/notes/notes-number-systems-summary.tex
% =========================================================


% ---------------------------------------------------------
\subsection*{\texorpdfstring{The Natural Numbers $\mathbb{N}$}{The Natural Numbers N}}

\begin{axiombox}{One-Based Peano Axioms for $\mathbb{N}$}
Let $\mathbb{N}$ be a set with distinguished element $1$
and successor operation $n \mapsto n{+\!\!+}$.

\medskip
\begin{tabular}{@{} l p{0.80\linewidth} @{}}
\textbf{P1} & $1 \in \mathbb{N}$. \\[4pt]
\textbf{P2} & $n \in \mathbb{N} \Rightarrow n{+\!\!+} \in \mathbb{N}$.
             \hfill\textit{(successor closure)} \\[4pt]
\textbf{P3} & $n{+\!\!+} \neq 1$ for all $n \in \mathbb{N}$.
             \hfill\textit{(the first element has no predecessor)} \\[4pt]
\textbf{P4} & $n{+\!\!+} = m{+\!\!+} \Rightarrow n = m$.
             \hfill\textit{(successor is injective)} \\[4pt]
\textbf{P5} & If $P(0)$ holds and $P(n) \Rightarrow P(n{+\!\!+})$,
             then $P(n)$ holds for all $n \in \mathbb{N}$.
             \hfill\textit{(induction)} \\
\end{tabular}
\end{axiombox}

\begin{remark*}
P1--P4 rule out wrap-around, ceiling, and rogue elements.
P5 eliminates anything unreachable from $1$ by finite iteration of the
successor.
Addition, multiplication, and order are not axioms -- they are \emph{defined}
recursively from the successor and then shown to satisfy familiar laws by
induction.
\end{remark*}

\begin{remark*}[Well-ordering and induction are equivalent over $\mathbb{N}$]
The following are logically equivalent:
\begin{itemize}
  \item Principle of Mathematical Induction (P5).
  \item Well-Ordering Principle: every nonempty $A \subseteq \mathbb{N}$
        has a least element.
  \item Strong Induction: $(\forall k < n,\, P(k)) \Rightarrow P(n)$
        for all $n$ implies $P$ holds everywhere.
\end{itemize}
One must assume at least one; the others follow.
Well-ordering is the engine behind the Division Algorithm, B\'{e}zout's Identity,
and unique prime factorization.
\end{remark*}

% ---------------------------------------------------------
\subsection*{\texorpdfstring{The Whole Numbers $\mathbb{W}$}{The Whole Numbers W}}

\begin{axiombox}{$\mathbb{W}$ as $\mathbb{N}$ with an Adjoined Zero}
The whole numbers are
\[
\mathbb{W}=\mathbb{N}\cup\{0\},
\qquad 0\notin\mathbb{N}.
\]
The successor, addition, multiplication, and order on $\mathbb{W}$ extend the
corresponding structures on $\mathbb{N}$:
\begin{itemize}
  \item $S_{\mathbb{W}}(0)=1$, and $S_{\mathbb{W}}$ agrees with
        $S_{\mathbb{N}}$ on $\mathbb{N}$;
  \item $0$ is the additive identity;
  \item $0$ is absorbing for multiplication;
  \item $0$ is the least element of $\mathbb{W}$.
\end{itemize}
\end{axiombox}

\begin{remark*}[What $\mathbb{W}$ gains over $\mathbb{N}$, and what it lacks]
\textbf{Gain:} an additive identity. This makes
$(\mathbb{W},+,0)$ a commutative monoid and prepares the ordered-pair
construction of the integers.

\textbf{Lack:} additive inverses for nonzero elements. Subtraction is still
not an internal operation in general, so $\mathbb{W}$ is not a ring.
\end{remark*}

% ---------------------------------------------------------
\subsection*{\texorpdfstring{The Integers $\mathbb{Z}$}{The Integers Z}}

\begin{axiombox}{Axioms for $\mathbb{Z}$ --- Ordered Commutative Ring}

\medskip
\noindent\textbf{Addition}

\medskip
\begin{tabular}{@{} l p{0.78\linewidth} @{}}
\textbf{A1} & $a + b = b + a$ \hfill\textit{(commutativity)} \\[3pt]
\textbf{A2} & $(a + b) + c = a + (b + c)$ \hfill\textit{(associativity)} \\[3pt]
\textbf{A3} & $\exists\, 0 \in \mathbb{Z}$ such that $a + 0 = a$
             \hfill\textit{(additive identity)} \\[3pt]
\textbf{A4} & $\forall\, a,\; \exists\, {-a} \in \mathbb{Z}$ with
             $a + (-a) = 0$
             \hfill\textit{(additive inverses)} \\[6pt]
\end{tabular}

\noindent\textbf{Multiplication}

\medskip
\begin{tabular}{@{} l p{0.78\linewidth} @{}}
\textbf{M1} & $a \cdot b = b \cdot a$ \hfill\textit{(commutativity)} \\[3pt]
\textbf{M2} & $(a \cdot b) \cdot c = a \cdot (b \cdot c)$
             \hfill\textit{(associativity)} \\[3pt]
\textbf{M3} & $\exists\, 1 \in \mathbb{Z},\; 1 \neq 0,$ such that $a \cdot 1 = a$
             \hfill\textit{(multiplicative identity)} \\[6pt]
\end{tabular}

\noindent\textbf{Distributivity}

\medskip
\begin{tabular}{@{} l p{0.78\linewidth} @{}}
\textbf{D} & $a \cdot (b + c) = a \cdot b + a \cdot c$
            \hfill\textit{(distributivity)} \\[6pt]
\end{tabular}

\noindent\textbf{Order}

\medskip
\begin{tabular}{@{} l p{0.78\linewidth} @{}}
\textbf{O1} & $a \leq b$ or $b \leq a$ \hfill\textit{(totality)} \\[3pt]
\textbf{O2} & $a \leq b \land b \leq a \Rightarrow a = b$
             \hfill\textit{(antisymmetry)} \\[3pt]
\textbf{O3} & $a \leq b \land b \leq c \Rightarrow a \leq c$
             \hfill\textit{(transitivity)} \\[3pt]
\textbf{O4} & $b \leq c \Rightarrow a + b \leq a + c$
             \hfill\textit{(order + addition)} \\[3pt]
\textbf{O5} & $a \geq 0 \land b \geq 0 \Rightarrow ab \geq 0$
             \hfill\textit{(order + multiplication)} \\
\end{tabular}
\end{axiombox}

\begin{remark*}[What $\mathbb{Z}$ gains over $\mathbb{W}$, and what it lacks]
\textbf{Gain:} Axiom A4 -- additive inverses. Subtraction is always defined
in $\mathbb{Z}$; it is not in $\mathbb{W}$.

\textbf{Lack:} Multiplicative inverses. There is no $x \in \mathbb{Z}$
with $2x = 1$, so M4 (every nonzero element is invertible) fails.
In the language of abstract algebra, $\mathbb{Z}$ is a
\emph{commutative ordered ring} but not a field.
The passage to $\mathbb{Q}$ restores multiplicative inverses by formally
adjoining fractions.
\end{remark*}

% ---------------------------------------------------------
\subsection*{\texorpdfstring{Derived Properties of $\mathbb{Z}$ (from the Ring Axioms)}{Derived Properties of Z (from the Ring Axioms)}}

\begin{remark*}[Zero multiplication --- the step most often looked up]
The identity $a \cdot 0 = 0$ is not an axiom; it is derived.
The proof uses only A3, A4, and D:
\[
a \cdot 0
= a \cdot (0 + 0)
\overset{\text{D}}{=} a \cdot 0 + a \cdot 0.
\]
Adding $-(a \cdot 0)$ to both sides (A4) gives $0 = a \cdot 0$, i.e.\ $a \cdot 0 = 0$.
No order axioms, no well-ordering, no properties of $\mathbb{N}$ are needed ---
this holds in any ring.
\end{remark*}

% ---------------------------------------------------------
\subsection*{\texorpdfstring{Structural Comparison: $\mathbb{N}$, $\mathbb{W}$, $\mathbb{Z}$, $\mathbb{Q}$, $\mathbb{R}$}{Structural Comparison: N, W, Z, Q, R}}

\begin{tcolorbox}[colback=gray!6, colframe=gray!40, arc=2pt,
  left=6pt, right=6pt, top=4pt, bottom=4pt,
  title={\small\textbf{Axiom Comparison Table}},
  fonttitle=\small\bfseries]
\small
\renewcommand{\arraystretch}{1.3}
\begin{tabular}{@{} l l c c c c c @{}}
\toprule
\textbf{Group} & \textbf{Axiom} &
$\mathbb{N}$ & $\mathbb{W}$ & $\mathbb{Z}$ & $\mathbb{Q}$ & $\mathbb{R}$ \\
\midrule
Addition
  & A1 commutativity          & \checkmark & \checkmark & \checkmark & \checkmark & \checkmark \\
  & A2 associativity          & \checkmark & \checkmark & \checkmark & \checkmark & \checkmark \\
  & A3 identity ($0$)         & $\times$   & \checkmark & \checkmark & \checkmark & \checkmark \\
  & A4 inverses ($-a$)        & $\times$   & $\times$   & \checkmark & \checkmark & \checkmark \\
\midrule
Multiplication
  & M1 commutativity          & \checkmark & \checkmark & \checkmark & \checkmark & \checkmark \\
  & M2 associativity          & \checkmark & \checkmark & \checkmark & \checkmark & \checkmark \\
  & M3 identity ($1$)         & \checkmark & \checkmark & \checkmark & \checkmark & \checkmark \\
  & M4 inverses ($1/a$)       & $\times$   & $\times$   & $\times$   & \checkmark & \checkmark \\
\midrule
  & D distributivity          & \checkmark & \checkmark & \checkmark & \checkmark & \checkmark \\
\midrule
Order
  & O1--O3 total order        & \checkmark & \checkmark & \checkmark & \checkmark & \checkmark \\
  & O4--O5 order + arithmetic & \checkmark & \checkmark & \checkmark & \checkmark & \checkmark \\
  & Well-ordering             & \checkmark & \checkmark & $\times$   & $\times$   & $\times$   \\
\midrule
Completeness
  & Least upper bound         & $\times$   & $\times$   & $\times$   & $\times$   & \checkmark \\
\midrule
\multicolumn{2}{@{}l}{\textbf{Structure}} &
\textit{positive arithmetic} & \textit{comm.\ semiring} & \textit{comm.\ ring} &
\textit{ordered field} & \textit{complete ordered field} \\
\bottomrule
\end{tabular}

\medskip
\noindent\small $\checkmark$ = satisfied \quad $\times$ = fails
\end{tcolorbox}

\begin{remark*}[The extension hierarchy]
Each number system repairs exactly one deficiency of its predecessor:
\[
\mathbb{N}
\;\xrightarrow{+\;0}\;
\mathbb{W}
\;\xrightarrow{+\;\mathrm{A4}}\;
\mathbb{Z}
\;\xrightarrow{+\;\mathrm{M4}}\;
\mathbb{Q}
\;\xrightarrow{+\;\mathrm{LUB}}\;
\mathbb{R}.
\]
$\mathbb{N} \to \mathbb{W}$: adjoin $0$ so addition has an identity.
$\mathbb{W} \to \mathbb{Z}$: adjoin additive inverses so subtraction is
always defined.
$\mathbb{Z} \to \mathbb{Q}$: adjoin multiplicative inverses so division by
any nonzero element is always defined.
$\mathbb{Q} \to \mathbb{R}$: adjoin least upper bounds so every bounded
increasing sequence converges.

Note that $\mathbb{N}$ and $\mathbb{W}$ are the well-ordered systems in the
chain.
$\mathbb{Z}$ loses well-ordering the moment negative integers are introduced
(the set of all negative integers has no least element).
\end{remark*}

\begin{topicbox}{From Extension to Embedding}
\begin{exposition}
The hierarchy above describes what each new system adds. The next chapter
looks at a separate question: how the earlier systems are identified inside
the later ones. That identification is not merely a matter of writing subset
symbols. It is carried by structure-preserving embeddings that justify
treating a number from an earlier construction as the corresponding number in
a larger system.
\end{exposition}
\end{topicbox}
